\section{Continuous Assessment 1 (CA-1) --- Set A}
\label{sec:ca1_set_a}

\LPUPaperHeader{Continuous Assessment (CA-1)}{A}{60 Minutes}{30 Marks}{CO1 (Introduction \& Sustainability), CO2 (Natural Resources \& Ecosystems)}

\begin{center}
\sffamily\textbf{\color{AlertRed}Note: Attempt all six questions. Each question carries 5 marks.}
\end{center}

\questionitem{1}{
    Define the four primary spheres of the Earth system (Atmosphere, Hydrosphere, Lithosphere, Biosphere). Detail the vertical thermal stratification of the atmosphere, identifying the lapse rate in the troposphere and the mechanism responsible for thermal inversion in the stratosphere.
}{5}{CO1}{L2 Understand}

\questionitem{2}{
    Explain the Brundtland definition of Sustainable Development. Formulate the Ehrlich-Holdren IPAT equation ($I = P \times A \times T$) and calculate the target technological carbon intensity required if population expands by $25\%$ and affluence increases by $60\%$, while total emissions must be reduced by $30\%$.
}{5}{CO1}{L3 Apply}

\questionitem{3}{
    Distinguish between perpetual, renewable, and non-renewable natural resources with examples. State the exact renewability replenishment criterion and explain the consequence when harvest rate exceeds regeneration rate.
}{5}{CO2}{L2 Understand}

\questionitem{4}{
    State the Universal Soil Loss Equation (USLE: $A = R \cdot K \cdot LS \cdot C \cdot P$) and define each parameter. Discuss three biological and mechanical soil conservation techniques utilized in watershed management.
}{5}{CO2}{L3 Apply}

\questionitem{5}{
    Explain the Ghyben-Herzberg relation governing coastal aquifers. Calculate the saltwater wedge intrusion depth below sea level if over-pumping drops the freshwater water table from $+3.0\text{ m}$ to $+1.5\text{ m}$ above mean sea level.
}{5}{CO2}{L3 Apply}

\questionitem{6}{
    State Lindeman's 10\% Trophic Efficiency Rule. In a grassland ecosystem where solar insolation is $1.0 \times 10^7\text{ kcal/m}^2/\text{yr}$ and producers convert $1.0\%$ into NPP, compute the energy available to herbivores and primary carnivores. Why is the pyramid of energy always upright?
}{5}{CO2}{L4 Analyze}

\newpage
\subsection*{Solutions and Marking Rubrics --- CA-1 Set A}

\begin{solutionbox}[Solution to Question 2 --- IPAT Calculation]
\textbf{1. Brundtland Definition (1 Mark):} ``Development that meets the needs of the present without compromising the ability of future generations to meet their own needs'' (\textit{Our Common Future}, 1987).\\
\textbf{2. Formulation (1 Mark):} $I = P \times A \times T$, where $I = \text{Impact}$, $P = \text{Population}$, $A = \text{Affluence}$, $T = \text{Technology}$.\\
\textbf{3. Mathematical Substitution (2 Marks):}\\
Initial state: $I_0 = P_0 \times A_0 \times T_0$.\\
Future state: $P_1 = 1.25 P_0$, $A_1 = 1.60 A_0$, $I_1 = 0.70 I_0$.\\
$I_1 = P_1 \times A_1 \times T_1 \implies 0.70 (P_0 A_0 T_0) = (1.25 \times 1.60) (P_0 A_0) T_1 = 2.00 (P_0 A_0) T_1$.\\
$T_1 = \frac{0.70}{2.00} T_0 = 0.35 T_0$.\\
\textbf{4. Required Reduction (1 Mark):}\\
$\text{Reduction in } T = (1 - 0.35) \times 100\% = 65.0\%$. Technological intensity must decline by $65\%$.
\end{solutionbox}

\begin{solutionbox}[Solution to Question 5 --- Ghyben-Herzberg Coastal Intrusion]
\textbf{1. Principle \& Derivation (2 Marks):}\\
Under hydrostatic balance, a column of freshwater of height $(h + z)$ balances a column of saltwater of height $z$:
\[
\rho_f g (h + z) = \rho_s g z \implies z = \frac{\rho_f}{\rho_s - \rho_f} h = \frac{1.000}{1.025 - 1.000} h = 40 h
\]
\textbf{2. Numerical Calculations (2 Marks):}\\
Initial saltwater interface depth: $z_{\text{initial}} = 40 \times 3.0\text{ m} = 120.0\text{ m}$ below sea level.\\
Final saltwater interface depth: $z_{\text{final}} = 40 \times 1.5\text{ m} = 60.0\text{ m}$ below sea level.\\
\textbf{3. Upward Saline Intrusion (1 Mark):}\\
$\Delta z = 120.0 - 60.0 = 60.0\text{ m}$. The saltwater interface ascends by $60\text{ meters}$, severely contaminating coastal drinking water wells.
\end{solutionbox}

\begin{solutionbox}[Solution to Question 6 --- Lindeman Energetics]
\textbf{1. Lindeman's Law (1.5 Marks):}\\
Only $\approx 10\%$ of energy is transferred across consecutive trophic levels; $90\%$ is dissipated as metabolic respiration heat and unconsumed biomass.\\
\textbf{2. Energy Calculations (2 Marks):}\\
$\text{NPP (Producers, } E_1) = 1.0 \times 10^7 \times 0.010 = 100{,}000\text{ kcal/m}^2/\text{yr}$.\\
Herbivores (Primary Consumers, $E_2$): $E_2 = 100{,}000 \times 0.10 = 10{,}000\text{ kcal/m}^2/\text{yr}$.\\
Primary Carnivores (Secondary Consumers, $E_3$): $E_3 = 10{,}000 \times 0.10 = 1{,}000\text{ kcal/m}^2/\text{yr}$.\\
\textbf{3. Upright Invariance (1.5 Marks):}\\
The energy pyramid is strictly upright because energy transfer is unidirectional and governed by the Second Law of Thermodynamics (entropy increases, heat dissipates).
\end{solutionbox}
\newpage
